\documentclass[11pt]{article}
\usepackage[margin=1in]{geometry}
\usepackage{amsmath,amssymb,amsthm,booktabs}
\usepackage[hidelinks]{hyperref}
\newcommand{\per}{\operatorname{per}}
\newcommand{\E}{\mathbb E}
\newtheorem{theorem}{Theorem}
\newtheorem{lemma}[theorem]{Lemma}
\newtheorem{corollary}[theorem]{Corollary}
\title{Strong self-permanent contraction on every complex bipartite support}
\author{Leo Gao\\\small AI-assisted research note}
\date{November 26, 2026}
\begin{document}
\maketitle
\begin{abstract}
Let $C$ be a complex correlation matrix whose nonzero off-diagonal
entries form a bipartite graph. Put $S=\tfrac12\|C-I\|_F^2$.
We prove the strong self-Bapat--Sunder inequality
$\per(C\circ\overline C)\le\per C$ with no restriction on either
part size, rank, or phases. Equality holds only for disjoint rank-one
pairs and singleton blocks. In particular, the self-Chollet ratio is at
most $4^{-S/3}$. Our lower bound uses the occupation distribution of
two orthonormal families; our upper bound compares injections with walks
in a random map. A binary polynomial positivity check completes a
one-row estimate. This removes the width-five restriction from a
recent structured-support preprint. It does not settle unrestricted
Chollet; no definitive priority claim is made.
\end{abstract}

\section{The result}
Write $\per A=\sum_{\sigma}\prod_i a_{i,\sigma(i)}$.
Chollet~\cite{Chollet} asked whether
$\per(A\circ B)\le\per A\per B$ for all Hermitian positive
semidefinite $A,B$, where $\circ$ denotes entrywise product. By
Cauchy--Schwarz applied to permutation monomials, it suffices to prove
$\per(A\circ\overline A)\le(\per A)^2$ for each $A$.
This reduction concerns the unrestricted conjecture; knowledge for a
restricted class of $A$ proves the paired inequality when \emph{both}
factors are covered. For historical surveys and recent restricted
results see~\cite{Wanless,Pant,rankfour}. In particular, Pant and
Singh obtain a stronger multiplier inequality on the sign-restricted
real bipartite class; the arbitrary complex phases here need not
satisfy their termwise-positivity hypothesis.
Schubr~\cite{Schubr} proved the strong self-inequality for complex
bipartite supports for which one of the two sides has at most five
vertices (and additional nonbipartite classes). Our result retains the
strong inequality and allows both sides to have arbitrary size.
The unrestricted strong inequality is false~\cite{Drury}.

Let $C$ be an $n\times n$ correlation matrix, that is, $C\succeq0$
is Hermitian and has diagonal one. Its off-diagonal support is bipartite
if, after simultaneous row and column permutation,
\[
 C=C(X)=\begin{pmatrix} I_p& X\\ X^*& I_q\end{pmatrix},\qquad p+q=n.
\]
Positive semidefiniteness is equivalent to $\|X\|_{\rm op}\le1$
by the Schur complement. Set $s_i=\sum_{j=1}^q|x_{ij}|^2$ and
$S=\sum_i s_i$; each $s_i\in[0,1]$.

\begin{theorem}\label{main}
For every such $X$ one has
\[
 Q(C):=\per(C\circ\overline C)\ \le\ P(C):=\per C,
 \qquad P(C)\ge4^{S/3}.
\]
Equality $Q(C)=P(C)$ holds exactly when, up to permutation, $C$ is a
direct sum of singleton identity blocks and rank-one two-by-two
correlation matrices. In particular $Q(C)/P(C)^2\le4^{-S/3}$.
\end{theorem}

\begin{corollary}
If $A,B$ are positive semidefinite with bipartite off-diagonal support
(the two bipartitions may differ), then $\per(A\circ B)\le\per A\per B$.
For correlation matrices there is the additional factor
$4^{-(S_A+S_B)/6}$ on the right.
\end{corollary}
\begin{proof}
Cauchy--Schwarz gives
$|\per(A\circ B)|\le\sqrt{\per(A\circ\overline A)\per(B\circ\overline B)}$.
Use Theorem~\ref{main} in the normalized setting. For general matrices
with positive diagonals use diagonal congruences on the two factors
separately; if either has a zero diagonal entry it has a zero row, and
both sides vanish.
\end{proof}

\section{The Gram permanent as an occupation expectation}
Use the product standard complex Gaussian measure on $\mathbb C^d$ for
the polynomial inner product $\langle f,g\rangle=\E[\overline f g]$.
The monomials are orthogonal and have squared norm
$\|z^\alpha\|^2=\alpha!=\prod_k\alpha_k!$. For linear forms $L_j$
with Gram matrix $A=(\E\overline{L_i}L_j)$, Wick's pairing formula gives
$\|\prod_j L_j\|^2=\per A$; this also follows by comparing
coefficients in the complex Gaussian exponential moment.

Realize the forms for $C$ as $z_1,\ldots,z_p,L_1,\ldots,L_q$
in dimension at most $p+q$, extending the first orthonormal family to
an orthonormal basis. The $L_j$ are orthonormal, and
$L_j(z)=\sum_i u_{ij}z_i$ with $u_{ij}=x_{ij}$ for $i\le p$.
Put $P=\prod_{j=1}^qL_j=\sum_\alpha c_\alpha z^\alpha$.
Then $\|P\|^2=\per I_q=1$; hence
$\mathbb P(N=\alpha)=\alpha!|c_\alpha|^2$ is a probability distribution.
Adding one factor $z_i$ increments its monomial exponent, whence
\begin{equation}\label{fock}
 \per C=\|z_1\cdots z_p P\|^2=\E\prod_{i=1}^p(1+N_i).
\end{equation}

For $i\le p$ the squared norms of $\partial_i P$ and $\partial_i^2 P$
are $\E N_i$ and $\E[N_i(N_i-1)]$: differentiate monomials and use
their norms. Products of distinct subsets of the $L_j$ are orthonormal
for distinct subsets, since these forms are coordinates in another
orthonormal basis. Differentiating the factored expression therefore yields
\begin{align}
 \E N_i&=\sum_j|x_{ij}|^2=s_i,\nonumber\\
 \E N_i(N_i-1)&=4\sum_{j<k}|x_{ij}|^2 |x_{ik}|^2
              =2\left(s_i^2-\sum_j|x_{ij}|^4\right)\le2s_i^2.
 \label{mom}
\end{align}
These identities do not assert independence of the occupations.

For a nonnegative integer variable $N$ with mean $s>0$, its size-biased
law assigns mass $k\mathbb P(N=k)/s$ to $k$. The function
$g(u)=\log(1+u)/u=\int_0^1(1+tu)^{-1}\,dt$ (continuously extended
at zero) is decreasing and convex. Jensen under the size-biased law
and \eqref{mom} give
\[
 \E\log(1+N_i)\ge s_i\,g(\E N_i^2/s_i)
 \ge s_i g(1+2s_i)\ge s_i g(3)=s_i\log(4)/3 .
\]
The zero-mean case is immediate. Taking log in \eqref{fock} and using
ordinary Jensen proves $P(C)\ge4^{S/3}$. Keeping the exact fourth
moment $q_i=\sum_j|x_{ij}|^4$ in \eqref{mom} also yields
\begin{equation}\label{lb}
 \log P(C)\ge\sum_{i:s_i>0} s_i g(1+2s_i-2q_i/s_i).
\end{equation}

\section{The squared-entry permanent}
For every rectangular matrix $Y$, partitioning a permutation by the
sets $I,J$ of moved vertices on the two sides gives
\[
 \per C(Y)=\sum_{|I|=|J|}|\per Y[I,J]|^2 .
\]
Indeed each of the two cross directions is a bijection between the
selected sets; all remaining vertices must be fixed.
Apply this to $W=(|x_{ij}|^2)$. Put $G=WW^{\mathsf T}$, a real
nonnegative matrix. For each row set $I$, expansion of
$\per G[I,I]$ includes every contribution to
$\sum_{J:|J|=|I|}(\per W[I,J])^2$: for each permutation of $I$,
expand its product of row inner products, first allowing only distinct
intermediate columns, then all choices. Distinct choices give exactly
the terms where the two bijections to $J$ use the same column set.
Thus $Q(C)\le\sum_I\per G[I,I]=\per(I_p+G)$.
Moreover $G_{ii}=q_i$ and $\sum_k G_{ik}\le s_i$, since every
column of $W$ has sum at most one.

Divide row $i$ of $I+G$ by $1+q_i$ and put
$r_i=(s_i-q_i)/(1+q_i)$. The resulting matrix is
$I+\operatorname{diag}(r_i)T$, where $T$ is zero-diagonal,
nonnegative, and row substochastic: take $T_{ik}=G_{ik}/(s_i-q_i)$
for $k\ne i$ when its denominator is positive, and zero on a null
row. Independently for every vertex $i$, choose a successor $k$ with
probability $T_{ik}$, or choose a cemetery symbol with the remaining
probability. The permanent of this matrix is
\[
 \E\prod_{D}\left(1+\prod_{i\in D}r_i\right),
\]
where $D$ ranges over directed cycles present in the random map.
Indeed the expansion selects disjoint directed cycles, each occurring
with probability equal to the product of its specified transitions;
vertices outside such cycles are fixed in the corresponding permutation.
For a cycle of length $k\ge2$, generalized H\"older and monotonicity
of finite-dimensional sequence norms give
$1+\prod_D r_i\le\prod_{i\in D}(1+r_i^k)^{1/k}
\le\prod_{i\in D}(1+r_i^2)^{1/2}$.
The cycles are disjoint, so multiplication and restoration of the row
factors show
\begin{equation}\label{ub}
 Q(C)\le\prod_i\sqrt{(1+q_i)^2+(s_i-q_i)^2}.
\end{equation}

\section{The scalar comparison}
For $s=s_i>0$ put $t=q_i/s_i^2\in[0,1]$. We prove
\[
 \tfrac12\log A(s,t)\ \le\ s\,g(1+2s(1-t)),
 \qquad A=(1+s^2t)^2+s^2(1-st)^2 .
\]
First, $\log A(s,t)/s$ is strictly increasing for $0<s\le1$ at
fixed $t$. Here is a short integer certificate for the only algebraic
inequality needed. Write $A_s=\partial A/\partial s$ and
$D=((sA_s)^2-A(A-1)^2)/s^4$, a polynomial of bidegree $(8,6)$.
The coefficients of
\[
 (1+u)^8(1+v)^6 D\big(u/(1+u),v/(1+v)\big)
\]
in increasing powers of $u$ (rows) and $v$ (columns), obtained by
binomial expansion, are the following integers:
\[
\begin{pmatrix}
3&30&117&228&237&126&27\\
24&220&796&1464&1456&748&156\\
83&688&2299&4024&3913&2008&425\\
162&1194&3684&6260&6186&3330&752\\
195&1250&3587&6100&6389&3746&925\\
148&800&2180&3840&4412&2912&812\\
69&300&813&1520&1971&1524&491\\
18&58&172&340&514&490&184\\
2&4&16&32&58&76&28
\end{pmatrix}.
\]
Thus $D>0$ on the closed unit square: in the equivalent Bernstein form
$D=\sum_{i,j}d_{ij}s^i(1-s)^{8-i}t^j(1-t)^{6-j}$ at least one
nonzero term remains at each boundary point. Since $A_s$ is positive
for small positive $s$ and cannot vanish on $(0,1]$, we have
$sA_s>\sqrt A(A-1)$. Finally $\log A\le(A-1)/\sqrt A$ for
$A\ge1$ (use $2\sinh(x/2)\ge x$ at $x=\log A$).
It follows that $sA_s/A>\log A$, as claimed.

It remains to show $g(3-2t)\ge\tfrac12\log(2+2t^2)$.
For $0\le t\le1/2$ the left hand side is at least $\log(4)/3$ and
the right hand side at most $\log(5/2)/2$; the former is strictly
larger since $16>(5/2)^3$.
For $1/2\le t\le1$ use the tangent to the convex function $g$ at
argument $1$. This lower bound for $g(3-2t)$ is affine in $t$,
equals $1/2$ at $t=1/2$, and equals $\log2$ at $t=1$.
The function $\tfrac12\log(2+2t^2)$ is convex on this interval, lies
strictly below $1/2$ at the left endpoint ($5/2<e$), and equals
$\log2$ at the right. The affine bound therefore suffices and is
strict for $t<1$. Combining both steps,
\[
 \frac{\log A(s,t)}{2s}\le
 \frac{\log A(1,t)}2\le g(3-2t)
 \le g(1+2s(1-t)).
\]
Equality requires $s=t=1$. Combining \eqref{ub} and \eqref{lb} proves
$Q(C)\le P(C)$. Equality forces every nonzero row of $W$ to have
$s_i=t_i=1$, hence exactly one entry, equal to one. These entries lie
in distinct columns by the column sum bound. Conversely the resulting
blocks give equality directly.

\section{Scope and limits}
When $X$ is a unitary matrix of order $d$, the theorem applies to
the union of any two orthonormal bases and gives
$\per(C\circ\overline C)/(\per C)^2\le 4^{-d/3}$.
The Fourier matrices exhibit unbounded separation from scaled diagonal
dominance: $|X_{ij}|=d^{-1/2}$ and the off-diagonal absolute part of
$C$ has spectral radius $\sqrt d$. Nontrivial Fourier phases cannot be
removed by multiplying rows and columns by phases because the phases on
a four-cycle are invariant in product. The bounds are not sharp: for
$X=I_d$ both permanents are $2^d$. The lower bound uses orthonormality
on each side in two distinct places, the normalization of $P$ and the
orthogonality of distinct derivative summands. Neither property holds for
unrestricted positive semidefinite matrices. The result concerns the
self-conjugate pair, or two covered factors, not an arbitrary second
factor. We have not established priority for this special case.

\begin{thebibliography}{9}
\bibitem{Chollet} J. Chollet, Is there a permanental analogue to Oppenheim's inequality?,
\emph{Amer. Math. Monthly} 89 (1982), 57--58.
\href{https://doi.org/10.1080/00029890.1982.11995380}{doi:10.1080/00029890.1982.11995380}.
\bibitem{Wanless} G.-S. Cheon and I. M. Wanless, An update on Minc's survey of open problems
involving permanents, \emph{Linear Algebra Appl.} 403 (2005), 314--342.
\bibitem{Pant} P. Pant and R. Singh, Structural classes for Chollet's permanent conjecture,
\href{https://arxiv.org/abs/2604.24192}{arXiv:2604.24192v2} (2026).
\bibitem{Schubr} S. Schubr, Strong permanent inequalities on structured supports,
with an application to graph Laplacians, independent preprint (Nov. 21, 2026),
\href{https://doi.org/10.5281/zenodo.24409106}{doi:10.5281/zenodo.24409106}.
\bibitem{Drury} S. W. Drury, A counterexample to a question of Bapat and Sunder,
\emph{Electron. J. Linear Algebra} 31 (2016), 69--70.
\bibitem{rankfour} Y. Ma, Quantitative Chollet inequalities for matrices of rank at most four,
\href{https://arxiv.org/abs/2609.37461}{arXiv:2609.37461} (2026).
\end{thebibliography}
\end{document}
